\centering
\scriptsize
\setlength{\tabcolsep}{2pt}
\renewcommand{\arraystretch}{1.15}
\begin{tabular}{cccccc}
    \colorbox{green!17}{\strut magic} &
    \colorbox{green!17}{\strut version} &
    \colorbox{gray!12}{\strut string\_table\_size} &
    \colorbox{gray!12}{\strut globals\_count} &
    \colorbox{green!17}{\strut imports\_count} &
    \colorbox{gray!12}{\strut publics\_count}
\end{tabular}

\vspace{0.35em}

\begin{tabular}{cccc}
    \colorbox{gray!12}{\strut string\_table} &
    \colorbox{green!17}{\strut imports} &
    \colorbox{green!17}{\strut publics} &
    \colorbox{gray!12}{\strut code}
\end{tabular}

\vspace{0.5em}

\footnotesize
Зеленым отмечены поля и секции, добавленные или измененные по сравнению с исходным форматом

\vspace{0.8em}

\noindent
\begin{minipage}[t]{0.46\linewidth}
    \textbf{Новые записи}
    \begin{itemize}
        \item добавлены магическое число и версия формата байткода
        \item \texttt{imports}: импортированные модули
        \item \texttt{publics}: добавлен флаг для различия публичных функций и глобальных переменных
    \end{itemize}
\end{minipage}\hfill
\begin{minipage}[t]{0.5\linewidth}
    \textbf{Изменения в инструкциях}
    \begin{itemize}
        \item в \texttt{BEGIN\_CLOSURE} добавлен операнд с числом замкнутых переменных
        \item добавлена инструкция \texttt{FAIL\_KEEP}
        \item внешние ссылки кодируются отрицательными операндами через таблицу строк
        \item вместо инструкций \texttt{read}/\texttt{write}/\texttt{length}/\texttt{string} вызываются непосредственно функции среды исполнения
    \end{itemize}
\end{minipage}
